\documentclass[11pt,a4paper]{article}
\usepackage[T1]{fontenc}
\usepackage{lmodern,amsmath,amssymb,amsthm,microtype}
\usepackage[margin=27mm]{geometry}
\usepackage[colorlinks=true]{hyperref}
\newtheorem{proposition}{Proposition}
\newtheorem{corollary}[proposition]{Corollary}
\newcommand{\Z}{\mathbb Z}
\title{Rational formulas for points on a cyclic elliptic subgroup}
\author{}\date{10 November 2026}
\begin{document}\maketitle

\begin{proposition}
Let $E$ be an elliptic curve over a field $K$ of characteristic zero.
Suppose $P\in E(K)$ is nontorsion and $2E(K)=\Z P$.
Then $\operatorname{End}_K(E)=\Z$.
Every $K$-rational rational map $f:E^r\dashrightarrow E$ has the form
\[
                 f(Q_1,\ldots,Q_r)=Q_0+\sum_{i=1}^r[m_i]Q_i,
                 \qquad Q_0\in E(K),\quad m_i\in\Z.
\]
If its value belongs to $2E(K)$ at one tuple from $(\Z P)^r$,
then $Q_0=[m_0]P$ for an integer $m_0$, and
\[
 f([n_1]P,\ldots,[n_r]P)
                    =[m_0+m_1n_1+\cdots+m_rn_r]P
                    \qquad(n_i\in\Z).
\]
\end{proposition}
\begin{proof}
For $\phi\in\operatorname{End}_K(E)$, the point $\phi(P)$
belongs to $2E(K)$ because $P$ is a double. Hence
$\phi(P)=[m]P$ for some $m\in\Z$. The endomorphism
$\phi-[m]$ kills the nontorsion point $P$. A nonzero elliptic
endomorphism has finite kernel, so $\phi=[m]$.

The extension and rigidity theorem for maps to abelian varieties
implies that $f$ extends over $E^r$ and is a translation of a
homomorphism~\cite[Chapter I, Corollary 3.7]{Milne}.
A homomorphism from a product is the sum of its restrictions to
the individual factors. These restrictions are $K$-endomorphisms
of $E$, and therefore are integer multiplications by the first
paragraph. If one value on $(\Z P)^r$ is a double, subtracting
its integer-linear part puts $Q_0$ in $2E(K)$ as well.
\end{proof}

\begin{corollary}
No $K$-rational point formula on $E^2$ can send
$([m]P,[n]P)$ to $[mn]P$ for every pair of integers $m,n$.
No such point formula on $E$ can send $[n]P$ to $[n^2]P$
for every integer $n$.
\end{corollary}
\begin{proof}
The label of a value of either rational point formula is affine
linear in the input labels by the proposition. Neither $mn$ nor
$n^2$ is such a function on the indicated integer tuples.
\end{proof}

The proposition restricts substitutions that give rational maps from
a product of elliptic curves to the same elliptic curve. It does not
exclude algebraic correspondences, additional equations on the inputs,
or constructions on other varieties. In particular it is not a lower
bound for the number of variables in a general Diophantine definition.

\begin{thebibliography}{1}
\bibitem{Milne}
J.~S. Milne, \emph{Abelian Varieties}, version 2.00 (2008),
Chapter I, Theorem 3.2 and Corollary 3.7, printed page 19.
\url{https://www.jmilne.org/math/CourseNotes/AV.pdf}.
The cited statements were checked directly in the author's text.
\end{thebibliography}
\end{document}
